\begin{table}[htbp]
\centering
\begin{threeparttable}
\caption{CHFS 2011 Cross-Section: Descriptive Statistics}
\label{tab:chfs_2011_summary}
{\small
\begin{tabular}{lcccccc}
\toprule
 & (1) & (2) & (3) & (4) & (5) & (6) \\
 & $N$ & Mean & Std.\ Dev. & Min & Median & Max \\
\midrule
\multicolumn{7}{l}{\textit{Panel A: Expectations and Forecast Errors}} \\
\addlinespace
Inflation expectation (1--5) & 8,333 & 4.104 & 0.837 & 1 & 4.00 & 5.000 \\
High inflation exp.\ (binary) & 8,333 & 0.841 & 0.365 & 0 & 1.00 & 1.000 \\
Forecast error proxy & 8,333 & 0.104 & 0.837 & $-3$ & 0.00 & 1.000 \\
\addlinespace
\multicolumn{7}{l}{\textit{Panel B: Information Channels}} \\
\addlinespace
High-salience channels (0--3) & 8,438 & 1.437 & 0.785 & 0 & 1.00 & 3.000 \\
Low-salience channels (0--2) & 8,438 & 0.410 & 0.561 & 0 & 0.00 & 2.000 \\
Above-median media (binary) & 8,438 & 0.383 & 0.486 & 0 & 0.00 & 1.000 \\
\addlinespace
\multicolumn{7}{l}{\textit{Panel C: Controls}} \\
\addlinespace
Risk attitude (1--5) & 8,295 & 3.847 & 1.234 & 1 & 4.00 & 5.000 \\
Log(income $+ 1$) & 8,381 & 9.768 & 2.268 & 0 & 10.24 & 14.914 \\
\bottomrule
\end{tabular}
}
\begin{tablenotes}[flushleft]
\footnotesize
\item \textit{Notes:} Data are from the 2011 China Household Finance Survey. Inflation expectation is reverse-coded so higher values indicate more inflationary expectations. Forecast error proxy equals expected (reversed ordinal) minus realized ordinal (2012 CPI $\approx 2.6\%$).
\end{tablenotes}
\end{threeparttable}
\end{table}
